\documentclass[10pt]{article}
\usepackage[margin=1in]{geometry}
\usepackage{amsmath,amssymb}
\usepackage{booktabs}
\usepackage{microtype}
\usepackage[hidelinks]{hyperref}

\newcommand{\code}[1]{\texttt{#1}}

\title{Review of Rashid \& Nanzer, ``Frequency and Phase Synchronization in
Distributed Antenna Arrays Based on Consensus Averaging and Kalman
Filtering''\\[2pt]
\large IEEE Transactions on Wireless Communications, vol.~22, no.~4,
pp.~2789--2802, April 2023. DOI 10.1109/TWC.2022.3213788}
\author{Chandler Stevenson}
\date{July 28, 2026}

\begin{document}
\maketitle

\section{Summary}
The paper addresses joint frequency and phase alignment of $N$ oscillators
in an open-loop coherent distributed array---one with no feedback from the
beamforming destination. Each node locally broadcasts its oscillator
frequency and phase to its neighbors over an undirected connectivity graph.
In the baseline algorithm, \emph{decentralized frequency and phase consensus}
(DFPC), every node iteratively replaces its state with a weighted
(Metropolis--Hastings) average of its neighbors'; the doubly stochastic
mixing matrix $\mathbf{W}$ drives all nodes to the arithmetic mean of their
initial states, with convergence speed and residual error governed by the
second-largest eigenvalue modulus $\lambda_2$.

Errors are injected at each iteration from a practical statistical model:
oscillator frequency drift from a two-term Allan-deviation law
$\sigma_f = f_c\sqrt{\beta_1/T + \beta_2 T}$ (white FM plus random-walk FM,
$\beta_1{=}\beta_2{=}5\times10^{-19}$ for a quartz TCXO), intra-interval
drift, phase jitter drawn per interval from the oscillator's integrated
phase-noise power ($A = -53.46$\,dBc $\Rightarrow \sigma_\theta =
2.7$\,mrad), and frequency/phase estimation errors at their Cram\'er--Rao
bounds for an $L$-sample observation at a given SNR. The central analytical
result (their Eq.~27) is a closed-form steady-state total phase error for
DFPC on sparsely connected graphs: a root-sum-of-squares of five terms ---
oscillator drift over the update interval ($2\pi\sigma_f T$), intra-interval
drift, phase jitter, and the two estimation errors.

To reduce the residual at short update intervals, the paper integrates a
per-node two-state Kalman filter (state $[f_n,\theta_n]$, process noise from
the drift/jitter statistics, measurement noise at the CRLBs) whose posterior
means are then consensus-averaged (KF-DFPC). Monte Carlo results
($f_c = 1$\,GHz, $f_s = 10$\,MHz, $T = 0.1$\,ms typical, $10^3$ trials)
show: (i) KF-DFPC converges in dramatically fewer iterations than DFPC on
sparse graphs (e.g., 2 vs.~16 iterations at $N{=}100$, connectivity 0.05;
4 vs.~158 at $N{=}65$); (ii) it outperforms diffusion LMS, diffusion KF, and
the Kalman consensus information filter in various regimes; (iii) its
converged total phase error is essentially independent of SNR, unlike DFPC,
whose error at low SNR is estimation-dominated; and (iv) sweeping the update
interval reveals an optimum ($T \approx 20$\,ms at their parameters) where
drift-driven and estimation-driven errors balance---DFPC degrades on either
side, while KF-DFPC holds the theoretical floor for $T \lesssim 20$\,ms.
Throughout, an $18^\circ$ total phase error is used as the threshold for
retaining ${\geq}90\%$ of ideal coherent gain.

\section{Contributions}
\begin{enumerate}
\item A fully decentralized joint frequency-\emph{and}-phase consensus
algorithm (DFPC), extending prior decentralized work that aligned frequency
only, with convergence to the arithmetic average of initial states.
\item A closed-form steady-state total phase error for DFPC (Eq.~27),
derived by propagating per-iteration error injections through the consensus
recursion, with the graph dependence captured through $\lambda_2$.
\item A practical, parameterized error model connecting oscillator
data-sheet quantities (Allan deviation coefficients, integrated phase-noise
power) and estimator quality (CRLBs vs.\ $L$ and SNR) to array-level phase
coherence, including the effect of TDMA-shared observation time.
\item KF-DFPC: the integration of per-node Kalman filtering with consensus
averaging, including the covariance bookkeeping across the mixing step, at
per-node complexity $\mathcal{O}(\mathrm{card}\{\chi_n\}+8)$---no asymptotic
overhead relative to plain DFPC for dense graphs.
\item A thorough comparative simulation study (DFPC, KF-DFPC, DLMS, DKF,
KCIF) across array size, connectivity, SNR, and update interval,
establishing when Kalman filtering pays: short update intervals, low SNR,
and sparse connectivity.
\end{enumerate}

\section{Limitations}
The following are assessed from a full read of the published text.
\begin{enumerate}
\item \textbf{No propagation channel in the measurement model.} Nodes
observe each other's oscillator signal directly in AWGN (their Eq.~15).
All geometry- and channel-induced phases are lumped into a uniform constant
$\theta_n^e$ assumed corrected elsewhere. The identifiability problem this
hides---a one-way receiver observes oscillator \emph{plus} channel phase and
cannot separate them without reciprocity or calibration---is central to any
over-the-air realization and is not addressed. Multipath appears only as a
qualitative remark.
\item \textbf{Statistics-level validation.} The simulations inject Gaussian
draws from the assumed error distributions (Algorithm~1) rather than
processing waveforms through a receiver. The theory is therefore validated
against its own assumptions; estimator failures (missed detection, timing
error, cycle slips), quantization, and analog impairments are outside the
model.
\item \textbf{Instantaneous, idealized actuation.} Consensus updates are
applied as immediate state assignments each iteration. There is no
correction latency, no NCO quantization, and no control-loop dynamics; the
single mention of latency concerns CRLB observation windows. Consequently
the analysis cannot expose the error contribution of estimator uncertainty
propagated over an actuation delay, which dominates real closed-loop
implementations at short intervals.
\item \textbf{No closed form for the Kalman variant.} Eq.~27 covers DFPC
only; KF-DFPC's steady state is characterized purely by simulation against
the DFPC bound. The steady-state (Riccati) error floor of the filtered loop
is not derived.
\item \textbf{Oscillator model omissions.} The two-term Allan-deviation law
captures white FM and random-walk FM but not flicker FM; phase jitter is
drawn independently each interval rather than evolving as a continuous
process, so intra- and inter-interval phase noise are statistically
disconnected.
\item \textbf{Idealized coordination plane.} Time synchronization is assumed
complete; sample-clock offset and its accumulated timing drift are absent;
the local broadcasts of estimates are error-free and instantaneous.
\end{enumerate}

\section{Relevance to this project}
This is the closest published prior art to the error-floor analysis pursued
in this repository. Its Eq.~27 anchors the framing: two of its terms
(drift$\times$interval and estimation error) reproduce terms of our
decomposition, so any contribution must be stated as what Eq.~27 lacks ---
actuation latency as a parameter distinct from the update interval, coupled
to the Kalman steady-state posterior; an analytical flicker term; a
continuous phase-noise process; and validation by an impairment-complete
sampled-IQ simulation rather than draws from the assumed statistics. The
paper's $18^\circ$/90\%-gain threshold and its error-vs-update-interval
design chart (their Fig.~10) are the natural presentation conventions to
adopt and extend. Its consensus layer is complementary rather than
competing: it operates above an idealized pairwise exchange, which is
precisely the physical link our simulator models.

\end{document}
